% Generated by scripts/materialize_unified_paper.py; do not edit.
\section{Proof production, checking, and trust}\label{fin:sec:certification}

\subsection{From finite formulas to checked refutations}

The final proof-producing runs used CaDiCaL~\cite{Biere2024} at source commit
\[
\sha{c60730422e758ef1cebe7aeddf2dda31c996bf04}.
\]
An exit code \(20\) and the text \texttt{UNSATISFIABLE} were treated only as
requests to check the emitted proofs.  The final mathematical evidence is the
successful replay of the DRAT traces~\cite{Wetzler2014,Heule2016} by the
official \texttt{drat-trim} source at commit
\[
\sha{2e3b2dc0ecf938addbd779d42877b6ed69d9a985}.
\]
The principal recorded executable has SHA-256
\[
\sha{31d9e8e04bc76a65f0a18ea212ac44af4cf7bc398629495cec093ced772c50f4}.
\]
All three checks terminated with exit code zero and \texttt{s VERIFIED}.

\begin{table}[t]
\centering
\caption{Exact-six proof artifacts and replay scale.  Hashes are abbreviated
here and given in full in \cref{fin:tab:budget6-artifacts}.}\label{fin:tab:proof-scale}
\begin{tabular}{@{}crrrl@{}}
\toprule
branch & raw DRAT bytes & proof lines & compressed bytes & replay\\
\midrule
0 & 3493847713 & 6649939 & 656144450 & verified\\
1 & 2137671968 & 5590150 & 433552903 & verified\\
2 &  759008359 & 2837771 & 152784369 & verified\\
\bottomrule
\end{tabular}
\end{table}

For branch \(0\), a fresh replay in a temporary directory took \(273.448\)
seconds inside the checker and reported \(135152646\) resolution steps in
the retained core.  Branch \(1\) reported \(209897525\) resolution steps in
its independently checked run.  These timings are descriptive, not logical
premises.  The source commit and successful replay are the portable evidence;
the executable hash records the historical build and is not required to match
a fresh binary compiled on another platform.

\subsection{Semantic reconstruction}

The proof checker establishes unsatisfiability of the bytes in each DIMACS
file.  Separate semantic checkers establish that those bytes are the intended
relaxations.  For every branch they independently:

\begin{enumerate}[label=(\arabic*)]
\item parse the frozen seed and regenerate the lexicographic pair-variable
  map;
\item compare all \(161700\) triangle clauses in order;
\item regenerate the PySAT equality counter and compare its variables and
  clauses in order;
\item check the branch unit and the exact clause partition;
\item reconstruct every hitting clause as all \(153\) pair variables of one
  \(18\)-vertex mask;
\item verify the compressed and raw sizes and SHA-256 identities; and
\item verify the budget-five dependency and the three-edge cover.
\end{enumerate}

Branch \(0\) additionally reparses all seven cut sources and reproduces the
deduplicated union.  Branches \(1\) and \(2\) likewise reconstruct their
complete stored cut banks.  An earlier interrupted branch-1 trace fails gzip
integrity, remains quarantined, and is excluded from every manifest and
conclusion.

\subsection{Trust boundary}

The claim is proof-carrying, not formally verified in a proof assistant.  Its
earliest machine dependencies are:

\begin{itemize}
\item the bitset independent-set verifier used for the seed facts;
\item the PySAT implementation of the standard sequential counter;
\item the \texttt{drat-trim} proof kernel, together with its compiler and
  execution hardware; and
\item the scripts that map matrices and masks to DIMACS.
\end{itemize}

The audit narrows this boundary through independent matrix parsing, exhaustive
triangle enumeration, small-instance cross-checks of the clique routine,
counter truth-table tests, complete clause reconstruction, content hashes,
gzip integrity tests, and proof replay with a freshly compiled checker.  A
future LRAT conversion and formally verified LRAT checker could reduce the
trusted kernel further, but is not required for the conventional
proof-carrying SAT standard used here.
